\begin{table}[!t]
\centering
\footnotesize
\begin{tabular}{@{}l rr rrr@{}}
\toprule
 & \multicolumn{2}{c}{97 firms} & \multicolumn{3}{c}{1,500 simulated sites} \\
\cmidrule(lr){2-3} \cmidrule(lr){4-6}
Method & Correlation & Length ratio & RMSE & Inclusion & Length \\
\midrule
Normal (ebnm) & 0.974 & 0.937 & 0.730 & 88.8\% & 2.09 \\
Unimodal (ebnm) & 0.950 & 0.881 & 0.699 & 91.1\% & 1.89 \\
NPMLE (ebnm) & 0.984 & 0.873 & 0.688 & 77.5\% & 1.60 \\
bayesEfron $g$, plugged in & 0.999 & 0.996 & 0.688 & 89.7\% & 1.76 \\
bayesEfron & 1.000 & 1.000 & 0.688 & 89.9\% & 1.76 \\
\bottomrule
\end{tabular}
\caption{Comparison of bayesEfron with empirical Bayes fits to the same data using ebnm (version 1.1.42; locations estimated where applicable; 90\% equal-tailed intervals). For the 97 firms, the correlation of a method's posterior means with those of bayesEfron and the length of its intervals relative to bayesEfron; for the 1,500 simulated sites, the root mean squared error of the posterior means, the inclusion proportion and mean interval length. The plug-in row treats the posterior mean of $g$ from bayesEfron as known. The last two rows therefore differ only in accounting for uncertainty in $g$. The REML normal model of Section 3.2 gives similar results to the ebnm normal prior (inclusion 89.0\%, RMSE 0.730).}
\label{tab:ebnm}
\end{table}
